\documentclass[conference]{IEEEtran}

\usepackage[utf8]{inputenc}
\usepackage[T1]{fontenc}
\usepackage{amsmath,amssymb}
\usepackage{graphicx}
\usepackage{booktabs}

\setlength{\textfloatsep}{6pt}
\setlength{\floatsep}{4pt}
\setlength{\intextsep}{6pt}

\title{Semana 9 y 10 del Proyecto Tower Copter -- Grupo 22}

\author{
\IEEEauthorblockN{Nombre 1, Nombre 2, Nombre 3}
\IEEEauthorblockA{
Universidad de los Andes, Bogot\'a, Colombia\\
correo1@uniandes.edu.co, correo2@uniandes.edu.co, correo3@uniandes.edu.co}
}

\begin{document}

\maketitle

\begin{abstract}
Se presenta el desarrollo de las semanas 9 y 10 del proyecto Tower Copter para el grupo 22. Se obtiene la funci\'on de transferencia de la planta, se analiza su estabilidad y se ajusta el modelo con un polo adicional para trabajar con una din\'amica de segundo orden. Sobre esa planta se dise\~na un controlador PID en frecuencia y se verifican las especificaciones de la gu\'ia mediante MATLAB y Simulink.
\end{abstract}

\section{Semana 9: Modelado y An\'alisis}
La planta dada en la gu\'ia est\'a descrita por
\begin{equation}
2.4712\dot{y}(t)+y(t)=0.32029\,u(t).
\label{eq:edo}
\end{equation}
Aplicando transformada de Laplace con condiciones iniciales nulas,
\begin{equation}
2.4712\,sY(s)+Y(s)=0.32029\,U(s),
\end{equation}
y despejando $Y(s)/U(s)$, se obtiene
\begin{equation}
G_0(s)=\frac{Y(s)}{U(s)}=\frac{0.32029}{2.4712s+1}.
\label{eq:g0}
\end{equation}
Esta funci\'on de transferencia no tiene ceros y posee un polo en
\begin{equation}
p_1=-\frac{1}{2.4712}=-0.4047.
\end{equation}
Como el polo est\'a en el semiplano izquierdo, la planta es estable en lazo abierto. Adem\'as, al ser de primer orden, su respuesta al escal\'on es mon\'otona y no oscilatoria. La constante de tiempo es
\begin{equation}
\tau=2.4712~\text{s},
\end{equation}
por lo que el tiempo de establecimiento aproximado al $2\%$ es
\begin{equation}
t_{ss}\approx 4\tau=9.8848~\text{s}.
\end{equation}
Adem\'as, su ganancia DC es $G_0(0)=0.32029$, de modo que ante un escal\'on unitario la planta por s\'i sola no alcanza la referencia de valor $1$.

Debido a que en la semana 10 se piden especificaciones t\'ipicas de segundo orden, se ajust\'o la planta agregando un polo diez veces m\'as a la izquierda que el polo original. Esto da
\begin{equation}
p_2=10p_1=-4.0466,
\end{equation}
lo cual equivale a multiplicar por
\begin{equation}
\frac{1}{0.24712s+1}.
\end{equation}
La planta ajustada queda
\begin{equation}
G(s)=\frac{0.32029}{(2.4712s+1)(0.24712s+1)}.
\label{eq:g}
\end{equation}
Esta planta no tiene ceros y sus polos son
\begin{equation}
p_1=-0.4047, \qquad p_2=-4.0466.
\end{equation}
Ambos polos quedan sobre el eje real negativo del plano $s$, por lo que el sistema es estable en lazo abierto. Adem\'as, al ser dos polos reales negativos, la respuesta es sobreamortiguada y no presenta oscilaciones. En la simulaci\'on del escal\'on unitario se obtuvo un sobreimpulso nulo y un tiempo de establecimiento aproximado de $9.93$ s. Como el polo adicional est\'a mucho m\'as a la izquierda, el comportamiento dominante sigue asociado al polo original.

Para el dise\~no CAD preliminar se propone una estructura con contenci\'on completa de la h\'elice, topes mec\'anicos superior e inferior, materiales definidos y dimensiones del montaje. En la entrega en PDF deben verse los planos, las dimensiones principales, los materiales elegidos, las caracter\'isticas generales del montaje y la estrategia de seguridad. Esta estrategia consiste en impedir el contacto de la h\'elice con el usuario y limitar mec\'anicamente el recorrido. El CAD puede entregarse como anexo, tal como permite la gu\'ia.

\begin{figure}[!t]
\centering
\begin{minipage}{0.48\columnwidth}
\centering
\fbox{\parbox[c][1.7cm][c]{0.95\columnwidth}{\centering
Insertar aqu\'i\\
el pzmap}}
\end{minipage}\hfill
\begin{minipage}{0.48\columnwidth}
\centering
\fbox{\parbox[c][1.7cm][c]{0.95\columnwidth}{\centering
Insertar aqu\'i\\
la respuesta en lazo abierto}}
\end{minipage}
\caption{Resultados principales de la semana 9.}
\label{fig:sem9}
\end{figure}

\begin{figure}[!t]
\centering
\fbox{\parbox[c][1.8cm][c]{0.95\columnwidth}{\centering
Insertar aqu\'i el CAD preliminar con dimensiones, materiales y estrategia de seguridad}}
\caption{Espacio para el dise\~no CAD preliminar.}
\label{fig:cad}
\end{figure}

\section{Semana 10: Sintonizaci\'on PID}
Para el grupo 22 la gu\'ia exige
\begin{equation}
M_p<15\%, \qquad t_{ss}<16~\text{s}.
\label{eq:specs}
\end{equation}
La sintonizaci\'on se hizo en el dominio de la frecuencia usando la relaci\'on
\begin{equation}
\omega_c \approx \frac{6}{X}=\frac{6}{16}=0.375~\text{rad/s},
\end{equation}
y un margen de fase objetivo de $60^\circ$. Con \texttt{pidtune} se obtuvo el controlador
\begin{equation}
C(s)=K_p+\frac{K_i}{s}+K_ds
\end{equation}
con
\begin{equation}
K_p=2.2539,\quad K_i=1.4689,\quad K_d=0.7592.
\label{eq:gains}
\end{equation}
Se escogi\'o este m\'etodo porque la gu\'ia pide desempe\~no en t\'erminos de sobreimpulso y tiempo de establecimiento, y en clase se relacion\'o ese objetivo con la frecuencia de cruce y el margen de fase. Por eso resulta razonable ajustar el PID directamente en el dominio de la frecuencia.

El t\'ermino proporcional acelera la respuesta, el integral elimina el error en estado estacionario y el derivativo amortigua el transitorio. En tiempo, la se\~nal de control se interpreta como
\begin{equation}
u(t)=K_pe(t)+K_i\int_0^te(\tau)d\tau+K_d\frac{de(t)}{dt}.
\label{eq:u}
\end{equation}

Con el lazo cerrado
\begin{equation}
T(s)=\frac{C(s)G(s)}{1+C(s)G(s)},
\end{equation}
se obtuvieron las m\'etricas de la Tabla~\ref{tab:resumen}. La frecuencia de cruce lograda fue $0.3750$ rad/s y el margen de fase obtenido fue $73.70^\circ$. El sistema cumple la especificaci\'on porque el sobreimpulso es menor al $15\%$ y el tiempo de establecimiento es menor a $16$ s.

La salida y la se\~nal de control se presentar\'an en una misma gr\'afica comparativa, de manera que se vea c\'omo evoluciona la respuesta del sistema frente al esfuerzo de control. Adem\'as, se incluir\'a una captura de la implementaci\'on en Simulink. Comparando ambos lazos, en lazo abierto la planta es estable pero no sigue una referencia unitaria, mientras que en lazo cerrado la acci\'on integral del PID elimina el error estacionario y la respuesta cumple con $M_p<15\%$ y $t_{ss}<16$ s. La sinton\'ia obtenida es conservadora: logra un sobreimpulso bajo sin exigir una acci\'on de control excesivamente agresiva.

\begin{table}[!t]
\caption{Resumen de resultados del controlador}
\label{tab:resumen}
\centering
\scriptsize
\resizebox{\columnwidth}{!}{%
\begin{tabular}{ccccccccc}
\toprule
$K_p$ & $K_i$ & $K_d$ & $\omega_c$ & PM & $t_r$ & $t_p$ & $t_{ss}$ & $M_p$ \\
\midrule
2.2539 & 1.4689 & 0.7592 & 0.3750 & 73.70$^\circ$ & 4.52 s & 9.00 s & 12.48 s & 4.11\% \\
\bottomrule
\end{tabular}
\vspace{-2mm}
}
\end{table}

\begin{figure}[!t]
\centering
\fbox{\parbox[c][2.0cm][c]{0.95\columnwidth}{\centering
Insertar aqu\'i la gr\'afica comparativa con $y(t)$ y $u(t)$ superpuestas}}
\caption{Gr\'afica comparativa de la salida y la se\~nal de control.}
\label{fig:matlab}
\end{figure}

\begin{figure}[!t]
\centering
\begin{minipage}{0.48\columnwidth}
\centering
\fbox{\parbox[c][1.6cm][c]{0.95\columnwidth}{\centering
Insertar aqu\'i la captura\\
de la implementaci\'on en Simulink}}
\end{minipage}\hfill
\begin{minipage}{0.48\columnwidth}
\centering
\fbox{\parbox[c][1.6cm][c]{0.95\columnwidth}{\centering
Insertar aqu\'i otra captura\\
de resultados en Simulink}}
\end{minipage}
\caption{Resultados complementarios obtenidos en Simulink.}
\label{fig:simulink}
\end{figure}

\section{Conclusi\'on}
La planta original del Tower Copter es de primer orden y estable. Despu\'es de agregar un polo diez veces m\'as a la izquierda, se dise\~n\'o un PID en frecuencia que cumpli\'o para el grupo 22 con $t_{ss}=12.48$ s y $M_p=4.11\%$. El documento cubre los \'items pedidos por la gu\'ia: modelado por Laplace, polos y ceros, estabilidad, respuesta al escal\'on, dise\~no CAD, sintonizaci\'on PID, implementaci\'on en Simulink y an\'alisis de resultados.

\end{document}
